% ============================================================
\chapter{TỔNG QUAN VÀ CƠ SỞ LÝ THUYẾT}
\label{ch:tong_quan}
% ============================================================

Chương này trình bày nền tảng lý thuyết của báo cáo theo bốn mạch: (i) bài toán phân
đoạn ngữ nghĩa và đặc thù của ảnh UAV; (ii) các kiến trúc mạng encoder--decoder; (iii)
bài toán mất cân bằng lớp cùng họ các hàm mất mát xử lý nó; và (iv) các độ đo đánh
giá. Đây là cơ sở để diễn giải thiết kế thí nghiệm ở Chương~\ref{ch:du_lieu} và kết
quả ở Chương~\ref{ch:ket_qua}.

% ============================================================
\section{Phân đoạn ngữ nghĩa trong thị giác máy tính}
% ============================================================

\subsection{Vị trí bài toán}
Trong thị giác máy tính, các bài toán nhận dạng có thể xếp theo mức độ chi tiết của
đầu ra (Bảng~\ref{tab:cv_tasks}). \textbf{Phân đoạn ngữ nghĩa} (semantic segmentation)
nằm ở mức điểm ảnh: gán cho \emph{mỗi} điểm ảnh một nhãn lớp, nhưng không phân biệt
các thực thể cùng lớp (khác với phân đoạn thực thể --- instance segmentation).

\begin{table}[H]
\centering
\caption{Phân cấp các bài toán nhận dạng theo độ chi tiết đầu ra.}
\label{tab:cv_tasks}
\small
\begin{tabular}{l l l}
\toprule
\textbf{Bài toán} & \textbf{Đơn vị đầu ra} & \textbf{Ví dụ} \\
\midrule
Phân loại ảnh        & toàn ảnh           & ``ảnh này là rừng'' \\
Phát hiện đối tượng  & hộp bao            & ``cây ở vị trí (x,y,w,h)'' \\
\textbf{Phân đoạn ngữ nghĩa} & \textbf{từng điểm ảnh} & ``điểm ảnh này là cây lá kim'' \\
Phân đoạn thực thể   & điểm ảnh + thực thể & ``cây số 1, cây số 2\dots'' \\
\bottomrule
\end{tabular}
\end{table}

\subsection{Phát biểu hình thức}
Cho ảnh đầu vào $\mathbf{x}\in\mathbb{R}^{H\times W\times 3}$, mô hình học một ánh xạ
$f_\theta:\mathbb{R}^{H\times W\times 3}\rightarrow \mathbb{R}^{H\times W\times C}$ trả
về điểm số (logit) $z_{i,c}$ cho mỗi điểm ảnh $i$ và lớp $c$. Xác suất dự đoán thu
được qua hàm softmax:
\begin{equation}
p_{i,c} = \frac{\exp(z_{i,c})}{\sum_{k=1}^{C}\exp(z_{i,k})},
\qquad \hat{y}_i = \arg\max_c\, p_{i,c}.
\end{equation}
Tham số $\theta$ được tối ưu bằng cách cực tiểu hóa một hàm mất mát $\mathcal{L}$ trên
tập huấn luyện thông qua lan truyền ngược và hạ gradient.

\subsection{Tiến trình phát triển}
Cột mốc khởi đầu là mạng tích chập đầy đủ (Fully Convolutional Network ---
FCN)~\cite{long2015fcn}, thay các tầng kết nối đầy đủ bằng tích chập để giữ thông tin
không gian và cho phép dự đoán dày đặc. Từ FCN, ba dòng kiến trúc lần lượt nổi lên:
(1) \emph{encoder--decoder} với skip-connection (U-Net~\cite{ronneberger2015unet});
(2) \emph{tích chập giãn} thu ngữ cảnh đa tỉ lệ (DeepLab~\cite{chen2018deeplabv3plus});
và (3) \emph{cơ chế tự chú ý} của Transformer thị giác
(SegFormer~\cite{xie2021segformer}, Swin~\cite{liu2021swin}). Báo cáo này dùng dòng
encoder--decoder với skip-connection làm nền.

% ============================================================
\section{Phân đoạn ngữ nghĩa ảnh UAV và viễn thám}
% ============================================================

Ảnh thu từ UAV có một số đặc thù khiến bài toán khó hơn ảnh đời thường: cảnh rộng và
phức tạp, đối tượng đa tỉ lệ (cây lớn cạnh xe nhỏ), biến thiên mạnh theo góc chụp và
độ cao bay, cùng phân bố lớp rất lệch. Các khảo sát gần đây về phân đoạn ảnh UAV viễn
thám~\cite{uavreview2024} ghi nhận encoder--decoder dựa trên FCN vẫn là hướng chủ đạo,
trong khi các mô hình Transformer (ví dụ SegFormer cho ảnh UAV~\cite{xie2021segformer})
và biến thể DeepLabV3+~\cite{chen2018deeplabv3plus} ngày càng phổ biến nhờ khả năng
nắm ngữ cảnh tầm xa.

Riêng với môi trường rừng, bộ dữ liệu \textbf{Forest Inspection}~\cite{forestdataset}
--- được dùng trong báo cáo --- cung cấp ảnh UAV tổng hợp (AirSim) ở nhiều góc chụp và
độ cao, với nhãn ở mức điểm ảnh cho 11 lớp. Đặc điểm nổi bật là mất cân bằng lớp
nghiêm trọng và thay đổi miền theo điều kiện bay --- chính là hai thách thức trung tâm
mà báo cáo hướng tới (Mục~\ref{sec:imbalance}).

Về các công trình liên quan, hướng phân đoạn/giám sát rừng từ ảnh trên không khá đa
dạng. Trong bài toán \textbf{phát hiện cây đổ/gỗ chết} --- vốn là lớp hiếm và khó ---
các nghiên cứu gần đây dùng học sâu trên ảnh UAV để lập bản đồ gỗ chết trong rừng
phương bắc~\cite{deadwood2025}, hoặc phân đoạn thực thể từng thân cây đổ bằng mạng FCN
kết hợp đường biên chủ động trên ảnh cận hồng ngoại~\cite{fallentrees2021}. Hướng
\textbf{phân đoạn theo loài/lớp phủ} được khai thác qua chuỗi ảnh trên không theo thời
gian~\cite{treeseg2024}, trong khi các bộ dữ liệu chuyên biệt cho học ngữ nghĩa rừng
tiếp tục được công bố~\cite{forestsemantic2024}. Điểm chung của các công trình này là
\textbf{mất cân bằng lớp} và \textbf{biến thiên điều kiện thu nhận} --- đúng hai thách
thức mà báo cáo tập trung giải quyết ở cấp hàm mất mát.

% ============================================================
\section{Kiến trúc Encoder--Decoder}
% ============================================================

\subsection{Ý tưởng chung}
Phần \textbf{encoder} (backbone) trích đặc trưng và giảm dần độ phân giải không gian
(tăng trường tiếp nhận, tăng số kênh ngữ nghĩa); phần \textbf{decoder} (head) khôi phục
độ phân giải để dự đoán nhãn ở mức điểm ảnh. Thông tin chi tiết bị mất khi hạ mẫu được
khôi phục nhờ \textbf{skip-connection} nối đặc trưng nông từ encoder sang decoder
(Hình~\ref{fig:encdec}). Đây là khuôn mẫu chung của hầu hết mô hình hiện đại.

\begin{figure}[H]
\centering
\resizebox{0.95\textwidth}{!}{%
\begin{tikzpicture}[
  font=\footnotesize, >=stealth,
  blk/.style={draw, rounded corners=2pt, minimum height=0.72cm, align=center},
  enc/.style={blk, fill=blue!12},
  dec/.style={blk, fill=orange!20},
  bott/.style={blk, fill=green!18},
  io/.style={blk, fill=gray!12}]
  % Input / Output
  \node[io, minimum width=1.7cm] (in)  at (-2.4,2.4) {Ảnh vào\\$H{\times}W{\times}3$};
  \node[io, minimum width=1.7cm] (out) at ( 7.0,2.4) {Bản đồ nhãn\\$H{\times}W{\times}C$};
  % Encoder (trái, đi xuống) - hạ mẫu
  \node[enc, minimum width=2.6cm] (e1) at (0,2.4)  {Enc-1};
  \node[enc, minimum width=2.1cm] (e2) at (0,1.2)  {Enc-2};
  \node[enc, minimum width=1.6cm] (e3) at (0,0.0)  {Enc-3};
  \node[enc, minimum width=1.1cm] (e4) at (0,-1.2) {Enc-4};
  % Bottleneck
  \node[bott, minimum width=1.2cm] (b) at (2.3,-2.3) {Bottleneck};
  % Decoder (phải, đi lên) - nâng mẫu
  \node[dec, minimum width=1.1cm] (d4) at (4.6,-1.2) {Dec-4};
  \node[dec, minimum width=1.6cm] (d3) at (4.6,0.0)  {Dec-3};
  \node[dec, minimum width=2.1cm] (d2) at (4.6,1.2)  {Dec-2};
  \node[dec, minimum width=2.6cm] (d1) at (4.6,2.4)  {Dec-1};
  % Luồng chính
  \draw[->] (in) -- (e1);
  \draw[->] (e1) -- (e2); \draw[->] (e2) -- (e3); \draw[->] (e3) -- (e4);
  \draw[->] (e4) -- (b);  \draw[->] (b) -- (d4);
  \draw[->] (d4) -- (d3); \draw[->] (d3) -- (d2); \draw[->] (d2) -- (d1);
  \draw[->] (d1) -- (out);
  % Skip-connection
  \draw[->, dashed, gray] (e1) -- node[above]{\scriptsize skip} (d1);
  \draw[->, dashed, gray] (e2) -- (d2);
  \draw[->, dashed, gray] (e3) -- (d3);
  \draw[->, dashed, gray] (e4) -- (d4);
  % Nhãn nhóm
  \node[below=0.15cm of e4, blue!50!black] {\textit{Encoder (hạ mẫu)}};
  \node[below=0.15cm of d4, orange!60!black] {\textit{Decoder (nâng mẫu)}};
\end{tikzpicture}}
\caption{Sơ đồ kiến trúc encoder--decoder kiểu U-Net: encoder hạ mẫu để trích đặc
trưng, decoder nâng mẫu để dự đoán nhãn; skip-connection (nét đứt) nối đặc trưng cùng
độ phân giải nhằm giữ chi tiết không gian.}
\label{fig:encdec}
\end{figure}

\subsection{U-Net và UNet++}
U-Net~\cite{ronneberger2015unet} có cấu trúc đối xứng hình chữ U: nhánh co
(contracting) đóng vai trò encoder, nhánh giãn (expanding) đóng vai trò decoder, nối
với nhau bằng skip-connection trực tiếp ở mỗi mức độ phân giải. Nhờ đó decoder vừa có
ngữ nghĩa mức cao vừa giữ được chi tiết không gian, rất phù hợp cho đối tượng nhỏ và
biên.

\textbf{UNet++}~\cite{zhou2018unetpp} cải tiến U-Net bằng \emph{lưới skip-connection
lồng nhau, dày đặc} (nested dense skip pathways): giữa encoder và decoder chèn thêm các
khối tích chập trung gian tạo thành nhiều đường dẫn ngắn dần. Mục đích là \emph{thu hẹp
khoảng cách ngữ nghĩa} giữa đặc trưng nông của encoder và đặc trưng sâu của decoder
trước khi hợp nhất --- điều mà skip-connection đơn của U-Net bỏ qua. UNet++ còn hỗ trợ
\emph{giám sát sâu} (deep supervision), đặt mất mát ở nhiều nhánh đầu ra. Trong thực
nghiệm, UNet++ thường cải thiện độ chính xác ở biên và vật thể nhỏ --- lý do nó được
chọn làm kiến trúc nền (Hình~\ref{fig:unetpp}).

\begin{figure}[H]
\centering
\resizebox{0.9\textwidth}{!}{%
\begin{tikzpicture}[
  >=stealth, font=\footnotesize,
  nd/.style={circle, draw, minimum size=0.85cm, inner sep=0pt, fill=blue!8},
  bb/.style={circle, draw, minimum size=0.85cm, inner sep=0pt, fill=blue!22},
  nout/.style={circle, draw, minimum size=0.85cm, inner sep=0pt, fill=green!22},
  down/.style={->, thick, blue!55!black},
  up/.style={->, thick, orange!75!black},
  skip/.style={->, dashed, gray}]
  \def\dx{2.0}\def\dy{1.5}
  % Hàng i=0 (X^{0,j})
  \node[bb]  (x00) at (0*\dx, 0*\dy) {$X^{0,0}$};
  \node[nd]  (x01) at (1*\dx, 0*\dy) {$X^{0,1}$};
  \node[nd]  (x02) at (2*\dx, 0*\dy) {$X^{0,2}$};
  \node[nd]  (x03) at (3*\dx, 0*\dy) {$X^{0,3}$};
  \node[nout] (x04) at (4*\dx, 0*\dy) {$X^{0,4}$};
  % i=1
  \node[bb] (x10) at (0*\dx, -1*\dy) {$X^{1,0}$};
  \node[nd] (x11) at (1*\dx, -1*\dy) {$X^{1,1}$};
  \node[nd] (x12) at (2*\dx, -1*\dy) {$X^{1,2}$};
  \node[nd] (x13) at (3*\dx, -1*\dy) {$X^{1,3}$};
  % i=2
  \node[bb] (x20) at (0*\dx, -2*\dy) {$X^{2,0}$};
  \node[nd] (x21) at (1*\dx, -2*\dy) {$X^{2,1}$};
  \node[nd] (x22) at (2*\dx, -2*\dy) {$X^{2,2}$};
  % i=3
  \node[bb] (x30) at (0*\dx, -3*\dy) {$X^{3,0}$};
  \node[nd] (x31) at (1*\dx, -3*\dy) {$X^{3,1}$};
  % i=4
  \node[bb] (x40) at (0*\dx, -4*\dy) {$X^{4,0}$};
  % Down-sampling (backbone encoder)
  \draw[down] (x00)--(x10); \draw[down] (x10)--(x20);
  \draw[down] (x20)--(x30); \draw[down] (x30)--(x40);
  % Up-sampling (chéo lên)
  \draw[up] (x10)--(x01); \draw[up] (x20)--(x11); \draw[up] (x30)--(x21); \draw[up] (x40)--(x31);
  \draw[up] (x11)--(x02); \draw[up] (x21)--(x12); \draw[up] (x31)--(x22);
  \draw[up] (x12)--(x03); \draw[up] (x22)--(x13);
  \draw[up] (x13)--(x04);
  % Dense skip (ngang)
  \draw[skip] (x00)--(x01); \draw[skip] (x01)--(x02); \draw[skip] (x02)--(x03); \draw[skip] (x03)--(x04);
  \draw[skip] (x10)--(x11); \draw[skip] (x11)--(x12); \draw[skip] (x12)--(x13);
  \draw[skip] (x20)--(x21); \draw[skip] (x21)--(x22);
  \draw[skip] (x30)--(x31);
  % Chú thích cạnh
  \node[anchor=west, blue!55!black]   at (4.7*\dx,-1.3*\dy) {$\downarrow$ hạ mẫu (encoder)};
  \node[anchor=west, orange!75!black] at (4.7*\dx,-1.9*\dy) {$\nearrow$ nâng mẫu};
  \node[anchor=west, gray]            at (4.7*\dx,-2.5*\dy) {-\,-\,$\rightarrow$ skip dày đặc};
  \node[anchor=west, green!45!black]  at (4.7*\dx,-3.1*\dy) {$X^{0,j}$: nhánh ra (giám sát sâu)};
\end{tikzpicture}}
\caption{Kiến trúc UNet++ với lưới skip-connection lồng nhau, dày đặc. Cột $X^{i,0}$ là
backbone encoder (hạ mẫu); các nút $X^{i,j}$ ($j{>}0$) nối qua các đường nâng mẫu và
skip dày đặc, thu hẹp khoảng cách ngữ nghĩa trước khi hợp nhất. Các nút $X^{0,j}$ cho
phép giám sát sâu (deep supervision).}
\label{fig:unetpp}
\end{figure}

\subsection{Encoder: ResNet50 và các họ backbone}
Báo cáo dùng encoder \textbf{ResNet50}~\cite{he2016resnet}: mạng 50 tầng với
\emph{khối phần dư} (residual block) $\mathbf{y}=\mathcal{F}(\mathbf{x})+\mathbf{x}$.
Kết nối tắt giúp gradient lan truyền tốt qua mạng sâu, tránh suy biến (degradation).
ResNet50 dùng \emph{khối nút cổ chai} (bottleneck: $1{\times}1 \rightarrow 3{\times}3
\rightarrow 1{\times}1$) để cân bằng độ sâu và chi phí tính toán.

Tổng quát, có hai họ encoder:
\begin{itemize}[leftmargin=1.4cm]
  \item \textbf{CNN} (ResNet, EfficientNet, ConvNeXt~\cite{liu2022convnext}): trường
  tiếp nhận cục bộ, hiệu quả và mạnh về kết cấu/biên cạnh.
  \item \textbf{Transformer} (MiT/SegFormer~\cite{xie2021segformer},
  Swin~\cite{liu2021swin}): dùng \emph{tự chú ý} (self-attention), trong đó mỗi vị trí
  ``nhìn'' được toàn bộ ảnh, nắm ngữ cảnh tầm xa tốt hơn --- có lợi cho cảnh UAV rộng,
  nhưng tốn bộ nhớ và cần tốc độ học nhỏ hơn.
\end{itemize}
Báo cáo cố định ResNet50 để cô lập ảnh hưởng của hàm mất mát; việc khảo sát các họ
encoder/decoder khác thuộc kế hoạch mở rộng (Mục~\ref{sec:multi_arch}).

\subsection{Các họ decoder}
\label{sec:decoders}
Decoder quyết định cách hợp nhất đặc trưng đa mức và phục hồi độ phân giải:
\begin{itemize}[leftmargin=1.4cm]
  \item \textbf{Skip-connection} (U-Net, UNet++): ghép đặc trưng nông--sâu; mạnh cho
  vật thể nhỏ và biên.
  \item \textbf{ASPP} (DeepLabV3+~\cite{chen2018deeplabv3plus}): tích chập giãn nhiều
  tỷ lệ giãn $\rightarrow$ ngữ cảnh đa tỉ lệ không làm mất độ phân giải.
  \item \textbf{MLP head} (SegFormer): hợp nhất 4 mức đặc trưng bằng lớp MLP đơn giản;
  hiệu quả khi đi cùng encoder Transformer.
  \item \textbf{UPerNet} (FPN + PPM): kim tự tháp đặc trưng kết hợp gộp kim tự tháp;
  tổng quát cho nhiều loại encoder.
\end{itemize}

\subsection{Học chuyển giao (Transfer Learning)}
Encoder được khởi tạo bằng trọng số đã huấn luyện trên ImageNet ($\sim$1{,}2 triệu
ảnh). Mô hình ``mang theo'' các đặc trưng tổng quát (biên, kết cấu, hình dạng), nhờ đó
hội tụ nhanh hơn, cần ít dữ liệu hơn và đạt độ chính xác cao hơn so với huấn luyện từ
đầu --- đặc biệt quan trọng với tập dữ liệu kích thước vừa như Forest Inspection. Trong
giai đoạn tinh chỉnh (fine-tuning), toàn bộ mạng được tối ưu trên dữ liệu mục tiêu.

% ============================================================
\section{Bài toán mất cân bằng lớp}
\label{sec:imbalance}
% ============================================================

\subsection{Định nghĩa và nguồn gốc}
Mất cân bằng lớp (class imbalance) xảy ra khi số điểm ảnh giữa các lớp chênh lệch lớn.
Trên dữ liệu rừng UAV, vài lớp nền (thảm thực vật, cây lá kim) chiếm đa số, trong khi
các lớp nhỏ/hiếm nhưng quan trọng (hàng rào, xe, cây đổ) chỉ chiếm phần rất nhỏ. Có thể
phân biệt hai dạng:
\begin{itemize}[leftmargin=1.4cm]
  \item \textbf{Mất cân bằng tần suất} (frequency imbalance): một số lớp áp đảo về số
  điểm ảnh.
  \item \textbf{Mất cân bằng độ khó} (hard-example imbalance): rất nhiều điểm ảnh
  ``dễ'' (nền) lấn át gradient so với số ít điểm ảnh ``khó'' (biên, vật thể nhỏ).
\end{itemize}

\subsection{Vì sao gây hại?}
Hàm mất mát trung bình theo điểm ảnh (như Cross-Entropy) cộng đóng góp của mọi điểm
ảnh như nhau. Khi một lớp chiếm $99\%$ điểm ảnh, gradient tổng bị nó chi phối; mô hình
có thể đạt độ chính xác điểm ảnh cao chỉ bằng cách ``bỏ qua'' lớp hiếm. Đây là
\textbf{nghịch lý độ chính xác} (accuracy paradox): chỉ số tổng đẹp nhưng các lớp thiểu
số --- thường là mục tiêu quan tâm --- lại bị dự đoán rất kém.

\subsection{Các nhóm giải pháp}
Giải pháp xử lý mất cân bằng chia thành ba mức:
\begin{enumerate}[leftmargin=1.4cm]
  \item \textbf{Mức dữ liệu:} lấy mẫu lại (oversampling lớp hiếm), tăng cường dữ liệu,
  copy-paste đối tượng hiếm.
  \item \textbf{Mức hàm mất mát:} đặt trọng số lớp, đổi dạng hàm (Focal, Dice,
  Tversky), hoặc kết hợp --- là trọng tâm của báo cáo.
  \item \textbf{Mức đánh giá:} báo cáo chỉ số theo từng lớp (per-class IoU) thay vì chỉ
  nhìn mIoU tổng.
\end{enumerate}

% ============================================================
\section{Các hàm mất mát}
\label{sec:loss}
% ============================================================

Theo các khảo sát~\cite{jadon2020survey,ma2021loss,losssurvey2023}, hàm mất mát cho
phân đoạn được chia thành ba nhóm: \emph{dựa trên phân phối} (CE, Focal), \emph{dựa
trên vùng} (Dice, Tversky, Lovász) và \emph{kết hợp} (compound). Phần này trình bày các
hàm dùng trong báo cáo và một số biến thể liên quan.

Ký hiệu: $N$ là số điểm ảnh hợp lệ, $C$ số lớp, $y_{i,c}\in\{0,1\}$ là nhãn one-hot,
$p_{i,c}$ là xác suất dự đoán, $p_{i,t}$ là xác suất của lớp đúng tại điểm $i$.

\subsection{Cross-Entropy (chuẩn đối sánh)}
Hàm mất mát cơ sở, thuộc nhóm dựa trên phân phối:
\begin{equation}
\mathcal{L}_{CE} = -\frac{1}{N}\sum_{i=1}^{N}\sum_{c=1}^{C} y_{i,c}\,\log p_{i,c}.
\label{eq:ce}
\end{equation}
Mọi điểm ảnh đóng góp như nhau nên với dữ liệu mất cân bằng, lớp hiếm gần như không ảnh
hưởng tới gradient tổng.

\subsection{Cross-Entropy có trọng số lớp}
Nhân mỗi lớp với trọng số $w_c$ tỷ lệ nghịch với tần suất để ``khuếch đại'' lớp hiếm:
\begin{equation}
\mathcal{L}_{wCE} = -\frac{1}{N}\sum_{i=1}^{N}\sum_{c=1}^{C} w_c\, y_{i,c}\,\log p_{i,c}.
\end{equation}
Báo cáo dùng công thức trọng số của ENet~\cite{paszke2016enet}:
\begin{equation}
w_c = \frac{1}{\ln(\kappa + f_c)}, \qquad \kappa = 1{,}02,
\label{eq:enet_w}
\end{equation}
với $f_c$ là tần suất điểm ảnh của lớp $c$; hằng số $\kappa$ chặn trọng số của lớp cực
hiếm không tăng vô hạn. \textbf{Hạn chế:} nếu nâng trọng số lớp hiếm quá mạnh, mô hình
dễ sinh dương tính giả và \emph{giảm} hiệu năng trên lớp phổ biến --- hiện tượng được
quan sát trong thực nghiệm (Chương~\ref{ch:ket_qua}).

\subsection{Focal Loss}
Focal Loss~\cite{lin2017focal} (đề xuất cho RetinaNet) xử lý mất cân bằng \emph{độ khó}
bằng cách hạ trọng số các ví dụ dễ:
\begin{equation}
\mathcal{L}_{Focal} = -\frac{1}{N}\sum_{i=1}^{N}
\alpha_t\,(1 - p_{i,t})^{\gamma}\,\log p_{i,t}.
\label{eq:focal}
\end{equation}
Thừa số điều biến $(1-p_{i,t})^{\gamma}$ làm đóng góp của ví dụ dễ ($p_{i,t}\!\to\!1$)
giảm mạnh, dồn gradient về các ví dụ khó. Với $\gamma = 0$, Focal trở về CE; báo cáo
dùng $\gamma = 2{,}0$ theo khuyến nghị gốc.

\subsection{Dice Loss}
Dice Loss~\cite{milletari2016vnet} thuộc nhóm dựa trên vùng, tối ưu trực tiếp mức độ
chồng lấp (xấp xỉ IoU) nên ít nhạy với mất cân bằng tần suất:
\begin{equation}
\mathcal{L}_{Dice} = 1 - \frac{1}{C'}\sum_{c=1}^{C'}
\frac{2\sum_i p_{i,c}\,g_{i,c} + \epsilon}{\sum_i p_{i,c} + \sum_i g_{i,c} + \epsilon},
\label{eq:dice}
\end{equation}
với $g_{i,c}$ là nhãn chuẩn, $\epsilon$ làm trơn tránh chia 0, $C'$ là số lớp được tính
(bỏ lớp nền \textit{Empty}). Bản tổng quát Generalised Dice~\cite{sudre2017gdl} nhân
thêm trọng số theo thể tích lớp.

\subsection{Một số biến thể dựa trên vùng}
Để hoàn chỉnh bức tranh, hai biến thể thường được nhắc tới:
\begin{itemize}[leftmargin=1.4cm]
  \item \textbf{Tversky Loss}~\cite{salehi2017tversky} tổng quát hóa Dice bằng cách
  đặt trọng số khác nhau cho dương tính giả (FP) và âm tính giả (FN):
  \begin{equation}
  \mathcal{L}_{Tversky} = 1 - \frac{TP}{TP + \alpha\,FP + \beta\,FN},
  \end{equation}
  cho phép ưu tiên \emph{recall} (giảm bỏ sót lớp hiếm) khi tăng $\beta$.
  \item \textbf{Lovász-Softmax}~\cite{berman2018lovasz} tối ưu trực tiếp chỉ số Jaccard
  (IoU) thông qua mở rộng lồi Lovász.
\end{itemize}
Các biến thể này không nằm trong phạm vi thí nghiệm hiện tại nhưng là hướng so sánh
tiềm năng (Chương~\ref{ch:ket_luan}).

\subsection{Hàm mất mát kết hợp (Compound)}
Mỗi hàm xử lý một khía cạnh; kết hợp chúng tận dụng ưu điểm bổ trợ: Focal lo điểm ảnh
khó, Dice lo chồng lấp vùng cho lớp nhỏ, CE giữ gradient ổn định. Báo cáo dùng:
\begin{equation}
\mathcal{L}_{combo} = \lambda_F\,\mathcal{L}_{Focal} + \lambda_D\,\mathcal{L}_{Dice}
+ \lambda_{CE}\,\mathcal{L}_{wCE},
\label{eq:combo}
\end{equation}
với bộ trọng số cố định $(\lambda_F,\lambda_D,\lambda_{CE}) = (1{,}0;\,1{,}0;\,0{,}5)$.
Việc chọn trọng số thủ công là điểm yếu --- dẫn tới ý tưởng học trọng số.

\subsection{Trọng số thích nghi (Uncertainty Weighting)}
\label{sec:adaptive}
Thay vì cố định $\lambda$, có thể \textbf{học} tỷ trọng giữa các thành phần theo cơ chế
bất định đa nhiệm của Kendall và cộng sự~\cite{kendall2018uncertainty}:
\begin{equation}
\mathcal{L}_{adaptive} = \sum_{k\in\{F,D,CE\}}
\left( \frac{1}{2\sigma_k^{2}}\,\mathcal{L}_k + \log\sigma_k \right),
\label{eq:adaptive}
\end{equation}
trong đó $\sigma_k$ (tham số học là $\log\sigma_k$) biểu diễn độ bất định của từng
thành phần và được tối ưu \emph{đồng thời} với tham số mô hình. Thành phần ``khó/nhiễu''
hơn tự nhận trọng số $\tfrac{1}{2\sigma_k^2}$ nhỏ hơn; số hạng $\log\sigma_k$ điều chuẩn
ngăn $\sigma_k$ tăng vô hạn.

\noindent\textit{Lưu ý:} do có số hạng $\log\sigma_k$ (có thể âm), giá trị
$\mathcal{L}_{adaptive}$ \textbf{có thể nhận giá trị âm} trong huấn luyện --- là hành vi
bình thường của công thức, không phải lỗi.

% ============================================================
\section{Độ đo đánh giá}
\label{sec:metrics}
% ============================================================

Mọi độ đo được suy ra từ \textbf{ma trận nhầm lẫn} (confusion matrix) $M\in\mathbb{R}^{C\times C}$,
trong đó $M_{ab}$ là số điểm ảnh có nhãn chuẩn $a$ nhưng bị dự đoán là $b$. Với lớp
$c$: $TP_c=M_{cc}$, $FP_c=\sum_{a\neq c}M_{ac}$, $FN_c=\sum_{b\neq c}M_{cb}$.

Chỉ số chính là \textbf{IoU} và trung bình của nó (mIoU):
\begin{equation}
\mathrm{IoU}_c = \frac{TP_c}{TP_c + FP_c + FN_c}, \qquad
\mathrm{mIoU} = \frac{1}{C'}\sum_{c=1}^{C'}\mathrm{IoU}_c.
\label{eq:miou}
\end{equation}
Bên cạnh đó là \textbf{Pixel Accuracy} (PA) và \textbf{mF1}:
\begin{equation}
\mathrm{PA} = \frac{\sum_c TP_c}{\text{tổng số điểm ảnh}}, \qquad
\mathrm{mF1} = \frac{1}{C'}\sum_{c=1}^{C'}\frac{2\,TP_c}{2\,TP_c + FP_c + FN_c}.
\end{equation}
Về hiệu quả tính toán dùng \textbf{Params} (triệu tham số), \textbf{FLOPs} (tỷ phép) và
\textbf{FPS} (khung/giây). Trong báo cáo, mIoU tính trên $C' = 10$ lớp (loại lớp nền
\textit{Empty}). \textbf{IoU theo từng lớp} đặc biệt quan trọng: nó bộc lộ hiệu năng
trên lớp hiếm mà mIoU tổng có thể che giấu --- đúng tinh thần xử lý mất cân bằng.

% ============================================================
\section{Tổng kết chương}
% ============================================================
Phân đoạn ngữ nghĩa ảnh UAV rừng là bài toán dự đoán nhãn mức điểm ảnh, giải bằng kiến
trúc encoder--decoder (báo cáo dùng UNet++/ResNet50, khởi tạo ImageNet). Thách thức cốt
lõi là mất cân bằng lớp, vốn khiến Cross-Entropy thiên vị lớp đa số. Các hàm mất mát
dựa trên vùng (Dice) và dựa trên độ khó (Focal), cùng cơ chế trọng số thích nghi, là
công cụ lý thuyết để khắc phục. Những nền tảng này định hình thiết kế thí nghiệm và
cách đánh giá trong các chương sau.
